\section{실패 자산화와 회귀 테스트}
탐색에서 찾은 실패는 한 번의 영상으로 끝나지 않고 재검증 가능한 테스트 자산이 되어야 한다.
모든 실행은 장면 파라미터, 시뮬레이션 모델, 작업 목표, 모델·소프트웨어 버전, 영상, 상태·접촉 기록,
기능 호출 결과를 위·변조 확인 정보와 함께 저장한다. 회귀 실행기는 이 기록을 읽어 동일
조건의 사례를 구성하고, 현재 버전으로 고정 예산 안에서 재실행한 뒤 원본과 비교한다.

\begin{enumerate}
\item 저장된 PASS/FAIL 사례를 읽고 동일 조건·장면의 사례를 묶는다.
\item 사례별 반복 수와 호출 상한을 정하고 입력·자산 호환성을 검사한다.
\item 현재 버전 조건을 기록한 뒤 재실행한다.
\item 원본과 신규 결과를 유지, 퇴행, 개선, 혼재 등으로 분류하고 영상·근거가 연결된 보고서를 생성한다(부록 E).
\end{enumerate}

정상 완료 사례도 함께 유지해 정책 업데이트 후 기존 성공이 깨지는지 확인한다. 본 과제에서 발견한
장애물 높이 경계(12/14\,cm)와 목표 점유 장면은 정책이나 현장 배치가 바뀔 때 우선 재실행할
회귀 테스트셋의 핵심 사례가 된다.
